\documentclass[11pt]{article}
\usepackage{amsmath,amssymb,amsthm,mathtools}
\usepackage{enumitem}
\usepackage[margin=1in]{geometry}
\setlength{\emergencystretch}{2em}

\newtheorem{theorem}{Theorem}
\newtheorem{definition}{Definition}
\newtheorem{nogo}{No-Go Theorem}
\newtheorem{remark}{Remark}

\title{Step 163: Scoped \(\Xi^{\rm BC}\) Adequacy Residual on \(\mathcal A_\eta/\mathcal K_\eta\)}
\author{Six Birds / RH Membrane Working Note}
\date{}

\begin{document}
\maketitle

\section{Inherited typed state}

The track is adequacy-oriented: smaller adequacy residual is better, and a
positive residual records missed adequacy rather than separation evidence.

The active Burnol/co-Poisson residual is
\[
  \Xi^{\rm BC}_{\ell,a}
  =
  \mathfrak B_{\ell,a}^{*}\Pi_{Y_a}\mathfrak B_{\ell,a}.
\]
The active operator is
\[
  C_\ell P_\eta,
  \qquad
  C_\ell=(I-P_\infty)M_{m_\ell}P_\infty,
  \qquad
  m_\ell(s)=e^{-\ell(1/2-s)}.
\]
The pulled-evaluator span is
\[
  H_\eta=
  \overline{\operatorname{span}\{J_a^*Y^a_{\rho,k}\}},
  \qquad
  P_\eta=\operatorname{Proj}_{H_\eta}.
\]
Step 162 fixed the Calkin context
\[
  \mathcal A_\eta=C^*(P_\infty,M_{m_\ell},P_\eta,I)\subseteq\mathcal B(H),
  \qquad
  \mathcal K_\eta=\mathcal A_\eta\cap\mathcal K(H),
\]
with quotient map
\[
  q_\eta:\mathcal A_\eta\to\mathcal A_\eta/\mathcal K_\eta.
\]
The completed quotient object is
\[
  q_\eta(C_\ell P_\eta)\in\mathcal A_\eta/\mathcal K_\eta.
\]

\section{Bridge gates}

\begin{center}
\begin{tabular}{lll}
\hline
Gate & Content & Status\\
\hline
G1 & \(\mathcal A_\eta/\mathcal K_\eta\) declared & framework\_defined\\
G2 & faithful symbol \(\sigma_B\) on \(q_\eta(C_\ell P_\eta)\) & open\_external\_proof\\
G3 & normal form \(C_\ell=U_\ell W_\ell+K_\ell\) & open\_external\_proof\\
G4 & lower-faithfulness of \(U_\ell\) on the relevant boundary range & open\_external\_proof\\
G5 & compact remainder \(K_\ell\in\mathcal K(H)\) & open\_external\_proof\\
\hline
\end{tabular}
\end{center}

\section{Scoped theorem}

\begin{theorem}[Scoped \(\Xi^{\rm BC}\) adequacy residual on \(\mathcal A_\eta/\mathcal K_\eta\)]
In the inherited Burnol/Sonine pulled-evaluator carrier and the typed Calkin
context \(\mathcal A_\eta/\mathcal K_\eta\), the following statements hold.
\begin{enumerate}[label=(\roman*)]
\item The operator
\[
  \Xi^{\rm BC}_{\ell,a}
  =
  \mathfrak B_{\ell,a}^{*}\Pi_{Y_a}\mathfrak B_{\ell,a}
\]
is the explicit adequacy residual carried by the current membrane packet.

\item The class \(q_\eta(C_\ell P_\eta)\) is the completed Calkin object whose
vanishing in \(\mathcal A_\eta/\mathcal K_\eta\) is the compact/tail-payability
gate for the active residual.

\item Under the split bridge theorem G2--G5, the vanishing of
\(q_\eta(C_\ell P_\eta)\) is equivalent to the boundary-Weyl compact/tail gate
\[
  W_\ell P_\eta\in\mathcal K(H),
\]
and hence to compact/tail-payability of
\[
  \Xi^{\rm BW}_{\ell,a}=P_\eta W_\ell^*W_\ell P_\eta
\]
in the bridge context.

\item Public-shadow non-promotion holds: Hardy/Hankel and hard-support symbol
computations do not decide \(q_\eta(C_\ell P_\eta)\) without an accepted
Calkin-faithful bridge into \(\mathcal A_\eta/\mathcal K_\eta\).

\item Finite-window Calkin blindness holds: finite singular-value and
finite-section evidence does not decide the completed class
\(q_\eta(C_\ell P_\eta)\).

\item The membrane packet at this carrier is diagnostic-complete: the active
residual is \(\Xi^{\rm BC}\), the current quotient object is
\(q_\eta(C_\ell P_\eta)\), and the only open external content at this leaf is
the split bridge theorem G2--G5.
\end{enumerate}
\end{theorem}

\begin{proof}
For (i), the residual is the Burnol/co-Poisson Schur-complement leftover
already fixed in the carrier record.  In the adequacy calculus, the residual is
positive and measures the component of the dissolving probe family not explained
by the native family.  The projection identity in the adequacy residual
calculus supplies the residual interpretation: after energy scaling, the
dissolving readout is projected away from the native explanation space.  Thus
the displayed \(\Xi^{\rm BC}_{\ell,a}\) is not a new operator introduced in this
step; it is the inherited explicit adequacy residual.

For (ii), steps 154 and 159 identify the pulled zero-evaluator residual kernel
with the action of \(C_\ell\) on the pulled-evaluator span \(H_\eta\).  Step 162
then declares the quotient in which this completed operator must be tested:
\(\mathcal A_\eta/\mathcal K_\eta\).  Therefore the active completed Calkin
object is \(q_\eta(C_\ell P_\eta)\), and vanishing of that class is exactly the
compact/tail-payability gate for the current residual.

For (iii), steps 160 and 161 prove only a conditional bridge equivalence.  If
G2 supplies a faithful boundary symbol on the declared quotient, G3 supplies
\(C_\ell=U_\ell W_\ell+K_\ell\), G4 supplies lower-faithfulness of \(U_\ell\) on
the boundary range relevant to \(H_\eta\), and G5 supplies the compact
remainder, then the compact part disappears in the quotient and
lower-faithfulness prevents \(U_\ell\) from killing the boundary class.  Hence
the Calkin class of \(C_\ell P_\eta\) vanishes exactly when the boundary-Weyl
gate \(W_\ell P_\eta\in\mathcal K(H)\) holds.  By step 159 this is the same
conditional reduction from \(\Xi^{\rm BC}\) to
\(\Xi^{\rm BW}_{\ell,a}=P_\eta W_\ell^*W_\ell P_\eta\).  No gate G2--G5 is
proved here.

For (iv), the public-shadow no-go from steps 160--162 applies.  A Hardy/Hankel
or hard-support computation lives in a different comparison context unless a
quotient-compatible bridge lands in \(\mathcal A_\eta/\mathcal K_\eta\).  Such
shadow data is non-injective on Calkin classes and cannot be promoted to a
statement about \(q_\eta(C_\ell P_\eta)\).

For (v), the finite-window Calkin blindness no-go from steps 158 and 160
applies.  A finite matrix corner or a finite list of singular values can be
held fixed while the tail is changed.  Since the completed Calkin class is a
tail statement, finite-window evidence remains moving-window support only until
a completed tail theorem is supplied.

For (vi), the residual tree at this carrier has been fully named:
\[
  \Xi^{\rm BC}
  \leadsto
  q_\eta(C_\ell P_\eta)
  \xrightarrow[\text{conditional}]{\text{G2--G5}}
  \Xi^{\rm BW},
\]
with public-shadow and finite-window shortcuts foreclosed and G2--G5 recorded
as precise remaining external proof obligations.  This is diagnostic-complete
in the sense of the construction cheat sheet: the residual is named, all bridge
statuses are audited, retained no-gos are in the ledger, and the remaining
external content is exactly stated.
\end{proof}

\section{Retained no-go records}

\begin{nogo}[Public-shadow non-promotion]
Hardy/Hankel and hard-support symbol computations cannot decide
\(q_\eta(C_\ell P_\eta)\) without an accepted bridge into
\(\mathcal A_\eta/\mathcal K_\eta\).
\end{nogo}

\noindent\textbf{Scope.} This forecloses routes that promote a public shadow
directly to a source statement on the Burnol/Sonine pulled-evaluator quotient.

\noindent\textbf{Proof status.} Framework-derived theorem-grade analytical
structural content inherited from steps 160--162.

\noindent\textbf{Nonclaim.} This does not rule out a true Calkin-faithful
bridge, a direct pulled-evaluator theorem, shifted co-Poisson factorization, or
a semilocal prolate theorem.

\begin{nogo}[Finite-window Calkin blindness]
Finite singular-value or finite-section evidence cannot decide the completed
Calkin class \(q_\eta(C_\ell P_\eta)\).
\end{nogo}

\noindent\textbf{Scope.} This forecloses finite-window evidence as a completed
Calkin certificate.

\noindent\textbf{Proof status.} Framework-derived theorem-grade analytical
structural content inherited from steps 158 and 160.

\noindent\textbf{Nonclaim.} This does not rule out a completed asymptotic
tail theorem, a genuine pulled-evaluator Weyl sequence, or a compact
decomposition on the declared carrier.

\section{Nonclaim boundary}

This theorem does not prove the Riemann hypothesis.  It does not prove
\(\Xi^{\rm BC}=0\).  It does not prove
\[
  C_\ell P_\eta\in\mathcal K(H).
\]
It does not prove
\[
  \|C_\ell P_\eta\|_{\rm e}>0.
\]
It does not accept any bridge source and does not close G2, G3, G4, or G5.  It
does not promote a public shadow or moving-window record to source status.  It
does not pivot to a new carrier.

\section{Verdict}

The final verdict of step 163 is
\[
  \boxed{\texttt{scoped\_residual\_theorem}.}
\]
The current carrier is diagnostic-complete with active residual
\(\Xi^{\rm BC}\), active quotient object \(q_\eta(C_\ell P_\eta)\), retained
no-gos, and open external content G2--G5.

\end{document}
